% IEEE conference-style technical report --- Generative AI Assignment 1.
% Build:  cd report && latexmk -pdf main.tex   (or upload the report/ folder to Overleaf)
% Every number in this report comes from an artifact in outputs/ (tables in report/tables are
% generated by the evaluation scripts). Remaining red TODOs: author details and the video link.
\documentclass[conference]{IEEEtran}
\usepackage[utf8]{inputenc}
\usepackage{amsmath,amssymb}
\usepackage{graphicx}
\usepackage{booktabs}
\usepackage{xcolor}
\usepackage{placeins}
\usepackage{hyperref}
\usepackage{cite}
\graphicspath{{figures/}{stitch/}}

% Figure/table includes that show a visible placeholder if an artifact is missing.
\newcommand{\artifactfig}[3][\columnwidth]{%
  \IfFileExists{figures/#2}{\includegraphics[width=#1]{#2}}{%
    \fbox{\parbox[c][3cm][c]{#1}{\centering\small\color{red} MISSING FIGURE: #2\\ (#3)}}}}
\newcommand{\stitchfig}[2][\columnwidth]{%
  \IfFileExists{stitch/#2}{\includegraphics[width=#1]{stitch/#2}}{%
    \fbox{\parbox[c][3cm][c]{#1}{\centering\small\color{red} MANUAL: Google Stitch evidence #2 not provided yet}}}}
% generated tables are typeset slightly smaller so that they fit one column
\newcommand{\artifacttable}[1]{\IfFileExists{tables/#1}{{\setlength{\tabcolsep}{2.5pt}\let\small\scriptsize\input{tables/#1}}}{%
  \begin{table}[t]\centering\fbox{\small\color{red} MISSING TABLE: tables/#1}\end{table}}}
\newcommand{\todo}[1]{\textcolor{red}{[TODO: #1]}}

\begin{document}

\title{Corruption-Aware Image Restoration with Autoencoders, Hard and Soft Expert Routing, and Style-Conditioned Face-to-Sketch Synthesis}

% Author details: name/roll number taken from the student's files; replace VIDEO_ID with the YouTube link.
\author{\IEEEauthorblockN{Muhammad Ali (i230008)}
\IEEEauthorblockA{Generative AI --- Assignment 1, Section B\\
\href{mailto:m.alenauman@gmail.com}{m.alenauman@gmail.com} --- Repository: \url{https://github.com/Ali-prog-spec/Restoration-Sketch-studio}\\
Demo video: \url{https://youtu.be/VIDEO_ID}}}

\maketitle

\begin{abstract}
We design, train, evaluate and deploy four generative systems. Tasks~1--3 restore Oxford-IIIT Pet images
corrupted at runtime by salt-and-pepper noise, Gaussian blur or rectangular occlusion: (1)~a single
convolutional denoising autoencoder with a genuine bottleneck, (2)~a corruption classifier that hard-routes
inputs to three specialist autoencoders (with an identity bypass for clean inputs), and (3)~a jointly
fine-tuned soft mixture-of-experts initialised from Task~2. Task~4 is a pix2pix-style conditional GAN that
translates FS2K face photographs into sketches of three styles through a learned style embedding used by
both generator and discriminator. All hyperparameters are tuned with Optuna, all runs are tracked with
MLflow, all deployed models are exported to ONNX and verified numerically, and the four systems are served
by one FastAPI + React application in Docker. On the official Oxford-IIIT Pet test set (36,690 corrupted
images), the mean SSIM rises from 0.672 for the corrupted inputs to 0.786 (universal autoencoder), 0.802
(hard routing, classifier accuracy 99.42\,\%) and 0.838 (soft mixture of experts). The face-to-sketch
generator reaches a test L1 of 0.111 and SSIM of 0.462, and changing only the style condition changes the
generated sketch by 0.10--0.12 mean absolute intensity.
\end{abstract}

\begin{IEEEkeywords}
denoising autoencoder, mixture of experts, routing, conditional GAN, pix2pix, Optuna, ONNX
\end{IEEEkeywords}

%=====================================================================
\section{Introduction}
Real images are often degraded in ways that are not known in advance: sensor noise produces isolated
black and white pixels, defocus produces blur, and occluders hide whole regions. A \emph{blind} restoration
model must remove whichever corruption is present without being told which one it is. This raises a
design question that we study experimentally: is it better to train one universal model, to first
recognise the corruption and send the image to a specialist (hard routing), or to let a learned gate mix
several specialists continuously (soft routing)? Hard routing is interpretable and lets every expert
specialise, but a single classification error sends an image to the wrong expert; soft routing is
differentiable and can hedge between experts when the input is ambiguous or contains more than one
corruption.

The second problem is paired image-to-image translation with a controllable output: converting a face
photograph into a sketch in one of three artist styles. Here the question is whether a single conditional
GAN can learn several styles through a learned style embedding, and whether that embedding genuinely
influences the output rather than serving as a label.

Our contributions are: (i)~a reproducible corruption pipeline with runtime training corruptions and
deterministic validation/test manifests; (ii)~a universal denoising autoencoder with a genuine bottleneck,
including an investigation of limited skip connections; (iii)~a hard-routed system evaluated with oracle
and predicted routing; (iv)~a soft mixture of experts initialised from Task~2 and fine-tuned jointly, with
a routing analysis; (v)~a style-conditioned pix2pix model on FS2K; and (vi)~one containerised web
application that serves all four systems from ONNX models.

\section{Related Work}
\textbf{Denoising autoencoders.} Vincent et al.~\cite{vincent2008dae,vincent2010sdae} introduced denoising
autoencoders that learn robust features by reconstructing clean inputs from corrupted ones. Convolutional
encoder--decoders with symmetric skip connections~\cite{mao2016rednet} and residual denoisers that predict
the noise~\cite{zhang2017dncnn} dominate restoration benchmarks, but their skip and residual paths carry
high-resolution input information around the bottleneck, which this assignment explicitly forbids. Zhao et
al.~\cite{zhao2017loss} showed that L2 correlates poorly with perceived quality and that mixtures of L1 and
(MS-)SSIM losses work better; SSIM is defined in~\cite{wang2004ssim} and its relation to PSNR is discussed
in~\cite{hore2010psnr}. Inpainting large holes is studied by context encoders~\cite{pathak2016context}.
Bilinear upsampling followed by convolution avoids the checkerboard artefacts of transposed
convolutions~\cite{odena2016checkerboard}.

\textbf{Mixture of experts.} Jacobs et al.~\cite{jacobs1991moe} proposed adaptive mixtures of local experts
with a softmax gate. Shazeer et al.~\cite{shazeer2017moe} scaled sparsely gated MoE layers and introduced an
importance (CV$^2$) loss against routing collapse; Switch Transformers~\cite{fedus2022switch} and
GShard~\cite{lepikhin2021gshard} use related load-balancing losses. Temperature-scaled softmax is used in
knowledge distillation~\cite{hinton2015distill}; confidence penalties based on entropy
regularise output distributions~\cite{pereyra2017confidence}.

\textbf{Conditional GANs.} GANs~\cite{goodfellow2014gan} and conditional GANs~\cite{mirza2014cgan} learn
generators through an adversarial game; DCGAN~\cite{radford2016dcgan} and~\cite{salimans2016improved}
introduced practical stabilisation techniques. Pix2pix~\cite{isola2017pix2pix} combines a U-Net
generator~\cite{ronneberger2015unet} with a PatchGAN discriminator and an L1 term. Conditions can be
injected by concatenating label maps~\cite{choi2018stargan}, by conditional normalisation~\cite{devries2017cbn,perez2018film}
or through a projection discriminator~\cite{miyato2018projection}. Image-generation quality is often
measured with LPIPS~\cite{zhang2018lpips} or FID~\cite{heusel2017fid}.

\textbf{Hyperparameter optimisation.} Optuna~\cite{akiba2019optuna} provides define-by-run search with the
TPE sampler~\cite{bergstra2011tpe} and early-stopping pruners related to successive
halving~\cite{jamieson2016sh,li2018hyperband}.

%=====================================================================
\FloatBarrier
\section{Dataset Preparation}
\subsection{Oxford-IIIT Pet (Tasks 1--3)}
The official trainval list (3,680 images)~\cite{parkhi2012pets} is split 80/20 with seed 42 into 2,944
training and 736 validation images; the 3,669 official test images are used only for the final evaluation.
Images are converted to RGB and resized to $128\times128$, and the same split is used by Tasks~1--3. Because
the official server limited each connection to about 150\,KB/s, the images were downloaded from a public
mirror and verified against the official source: the mirror's splits equal the official
\texttt{trainval.txt}/\texttt{test.txt}, and all 1,145 images recovered from a partial official download are
byte-identical (MD5).

\subsection{Corruption configuration}
\begin{table}[t]\centering\caption{Corruption configuration (assignment specification).}\label{tab:corruptions}\footnotesize
\setlength{\tabcolsep}{3pt}
\resizebox{\columnwidth}{!}{%
\begin{tabular}{lll}\toprule
Condition & Training (runtime, uniform) & Test levels (low / med / high)\\\midrule
Clean & original image & -- \\
Salt-and-pepper & $p\sim U(0.02,0.15)$, black/white 50/50 & $p=0.03 / 0.08 / 0.15$\\
Gaussian blur & $k\in\{3,5,7\}$, $\sigma\sim U(0.5,2.5)$ & $(3,0.7)/(5,1.5)/(7,2.5)$\\
Occlusion & 1--3 black rectangles, union 10--35\% & $\approx$10/20/35\%, 1/2/3 rect.\\\bottomrule
\end{tabular}}\end{table}
Table~\ref{tab:corruptions} lists the corruptions. During training a new condition and severity are sampled
every time an image is loaded, and a balanced batch sampler gives every batch exactly $B/4$ images of each
condition. Salt-and-pepper corrupts whole pixels; blur uses reflect padding; for occlusion the coverage is
the \emph{union} of the rectangles, so overlapping rectangles cannot under-cover. Validation (2,944 records,
four conditions per image) and test (36,690 records, clean plus three corruptions at three levels per image)
corruptions are stored once in deterministic manifests that hold the type, level, seed, probability, kernel,
$\sigma$, rectangle coordinates and coverage; a unit test confirms that reloading a manifest reproduces
identical images. The realised test occlusion coverage is 10.00\,\% (range 9.87--10.16\,\%), 19.83\,\%
(18.51--20.30\,\%) and 34.55\,\% (33.50--35.42\,\%) for the three levels.

\subsection{FS2K (Task 4)}\label{sec:fs2k_ref}
The official split (1,058 train / 1,046 test)~\cite{fan2022fs2k} is used; 15\,\% of the training portion
is held out for validation, stratified by style with seed 42, giving 899/159 images. Photos and sketches
are resized to $128\times128$; a paired augmentation (resize to 143, identical random crop and horizontal
flip) keeps every photograph aligned with its sketch, which a unit test verifies. The sketches are stored as
RGB but all three channels are identical, so the generator outputs one channel. Two properties of FS2K
matter for the analysis: the official test set is style-imbalanced (Style~1: 619, Style~2: 381,
Style~3: 46 images, whereas training is balanced), and style is confounded with the photo source (all
\emph{photo2} images are Style~3, while \emph{photo1} contains almost only Styles~1 and~2).

\section{Experimental Setup}
All models are implemented in PyTorch~2.14 (CUDA~12.6) and trained on one NVIDIA RTX~2060 (6\,GB) with an
Intel i7-9750H CPU, Python~3.11.9 and Optuna~5.0. Seed~42 is used for all splits, samplers and
initialisations. Tasks~1--3 use mixed-precision training; the GAN is trained in full precision for
stability. Every script supports SMOKE (tiny subset, used to test the whole pipeline), DEV and FINAL modes;
only FINAL results appear in this report. Every Optuna study uses the TPE sampler (seed~42), median pruning
and SQLite storage, and its first trial is the configuration suggested by the assignment, so each study
directly compares the suggested baseline with the tuned result. All runs (hyperparameters, per-epoch losses
and metrics, figures, checkpoints and the best Optuna configuration) are logged with MLflow.

For the restoration tasks the model-selection objective is
$J=0.5\,\mathrm{L1}+0.5\,(1-\mathrm{SSIM})$ on the validation manifest (lower is better). It combines
reconstruction error and structural similarity, as the assignment requires, but is deliberately independent
of the trainable loss weight $\alpha$: if trials were compared by their own training loss, Optuna could
lower the objective simply by changing $\alpha$. We report L1, SSIM and, as an additional standard
restoration metric, PSNR~\cite{hore2010psnr}, always together with the metrics of the corrupted input so
that the gain of restoration is visible.

%=====================================================================
\FloatBarrier
\section{Task 1: Universal Denoising Autoencoder}
\subsection{Methodology and architecture}
\begin{figure}[t]\centering\artifactfig{arch_autoencoder.pdf}{architecture}
\caption{Convolutional DAE. No skip connections: all information passes through the latent $z$.}\label{fig:ae}\end{figure}
The universal model (Fig.~\ref{fig:ae}) is a convolutional autoencoder without skip connections. Each
encoder stage halves the resolution with a stride-2 convolution followed by a second convolution
(BatchNorm~\cite{ioffe2015bn}, LeakyReLU); a $1\times1$ convolution then compresses the features into the
latent code $z$, and a mirrored decoder of bilinear upsampling and convolutions reconstructs the RGB image
through a sigmoid. Optuna chooses the bottleneck shape from five options ($8\times8\times\{16,32,64\}$,
$16\times16\times\{8,16\}$), i.e.\ 1,024--4,096 latent values for a 49,152-value input, a compression of
12--48$\times$. Every option is a genuine bottleneck, so no trial can improve by removing it. The selected
model uses three encoder stages with $C=48,96,192$ channels and a $16\times16\times8$ latent (2,048 values,
24$\times$ compression; 1.01\,M parameters). A unit test shows that, without skips, the decoder output
depends only on $z$.

To investigate limited skip connections we trained the same configuration with one skip at the coarsest
decoder level ($16\times16$), squeezed to 4 channels by a $1\times1$ convolution; there is no skip at higher
resolutions and no input--output residual.

\subsection{Loss}
$\mathcal{L}_{UDAE}=\alpha\,\mathcal{L}_{1}(x,\hat x)+(1-\alpha)(1-\mathrm{SSIM}(x,\hat x))$, with SSIM
($11\times11$ Gaussian window, $\sigma=1.5$) implemented in PyTorch and checked against scikit-image to
within $2\cdot10^{-3}$.

\subsection{Optuna search}
\artifacttable{optuna_task1.tex}
\begin{figure}[t]\centering\artifactfig{task1_optuna_history.png}{Optuna history}\\
\artifactfig{task1_optuna_importance.png}{Optuna importances}
\caption{Task 1 Optuna optimisation history and parameter importance.}\label{fig:t1opt}\end{figure}
The study ran 30 trials of 15 epochs (20 completed, 10 pruned). The assignment baseline (trial~0:
$\alpha=0.8$, learning rate $10^{-3}$, batch 32, $8\times8\times16$ bottleneck) reached $J=0.250$; the best
trial (29) reached $J=0.145$, a 42\,\% reduction. The fANOVA importances (Fig.~\ref{fig:t1opt}) attribute
most of the variance to dropout (0.62) and the learning rate (0.32): the best trials use almost no dropout
(0.006), because dropout on the latent code removes information the decoder needs. The selected $\alpha=0.53$
is lower than the suggested 0.8, i.e.\ structural similarity receives more weight. The selected
configuration was retrained for 80 epochs (best validation epoch 79: SSIM 0.798, PSNR 25.53\,dB).

\subsection{Results}
\artifacttable{task1_by_type.tex}
\artifacttable{task1_by_level.tex}
\begin{figure}[t]\centering\artifactfig{task1_training_curves.png}{training curves}
\caption{Task 1 training and validation curves.}\label{fig:t1curves}\end{figure}
\begin{figure}[t]\centering\artifactfig{task1_severity_ssim.png}{severity}
\caption{Task 1 test SSIM of the corrupted input and the restored output per corruption and severity.}\label{fig:t1sev}\end{figure}
\begin{figure*}[t]\centering
\artifactfig[0.48\textwidth]{task1_representative_1.png}{examples 1-6}\hfill
\artifactfig[0.48\textwidth]{task1_representative_2.png}{examples 7-12}
\caption{Twelve representative test examples (the image closest to the median $J$ of each condition and
level; three for clean): clean target, corrupted input, reconstruction and absolute error map.}\label{fig:t1rep}\end{figure*}
The training and validation curves (Fig.~\ref{fig:t1curves}) decrease together without a widening gap,
so the model does not overfit. On the test set the mean SSIM rises from 0.672 (inputs) to 0.786 and the
mean L1 falls from 0.050 to 0.039 (Tables~\ref{tab:task1_type}--\ref{tab:task1_level}). The benefit depends
strongly on the corruption. Salt-and-pepper is removed almost completely: SSIM rises from 0.382 to 0.825, and
the output quality hardly depends on the noise level (0.830/0.826/0.818 for $p=0.03/0.08/0.15$), because
isolated extreme pixels are easy to detect and the bottleneck cannot represent them. Occlusion is
improved moderately overall (0.713$\to$0.722): large holes gain most (0.546$\to$0.653 at 35\,\% coverage),
but at 10\,\% coverage the output is worse than the input (0.862$\to$0.781), because re-encoding the whole
image costs more detail than filling the small hole gains. For blur the model helps at medium and high
severity (0.806$\to$0.820, 0.692$\to$0.746); for mild blur (input SSIM 0.944) and clean inputs (output SSIM
0.831) the output is \emph{worse} than the input, because the 24$\times$ bottleneck cannot transmit fine
texture. Fig.~\ref{fig:t1sev} shows this crossover:
restoration gains grow with severity, while near-clean inputs lose detail. Averaged over corruptions, the
output SSIM is 0.813/0.792/0.739 for low/medium/high severity, while the input SSIM falls much faster
(0.803/0.625/0.481). The twelve examples in Fig.~\ref{fig:t1rep} show the same pattern: noise is removed and
edges are kept, occlusions are replaced by smooth colour that roughly continues the surroundings, and the
error maps concentrate at fine textures (fur, grass) and inside former occlusions.

\textbf{Limited skips.} The ablation with one narrow skip at $16\times16$ improves every condition (overall
SSIM 0.839 vs.\ 0.786; clean 0.898 vs.\ 0.831; occlusion 0.771 vs.\ 0.722). Even a 4-channel skip at a
coarse resolution carries enough spatial detail to sharpen the outputs. Because the assignment requires all
information to pass through a genuine bottleneck, the deployed model is the no-skip version; the ablation
quantifies the fidelity this constraint costs.

\subsection{Failure analysis}
\begin{figure}[t]\centering\artifactfig{task1_failures_1.png}{failure cases}
\caption{Task 1 failure cases: the worst output per corruption type and the largest degradations of a
clean input.}\label{fig:t1fail}\end{figure}
Fig.~\ref{fig:t1fail} shows seven automatically selected failures. (1)~A spotted cat on a black background
with 35\,\% occlusion (SSIM 0.723$\to$0.299): the model fills the black background with the grey-green
average colour it learned for pet backgrounds, so most of the error is outside the cat; the model has no
notion of which pixels are visible. (2)~A dog in dense flowers under strong blur (0.335$\to$0.364): the
texture is beyond what the bottleneck can recover, and the output is only slightly sharper. (3)~A beagle
in a meadow with 15\,\% salt-and-pepper (0.372$\to$0.504): the noise is removed but so is the grass texture.
(4--7)~Clean images with high-frequency patterns (striped blanket, grass, fence) lose the pattern
(SSIM 1.0$\to$0.54--0.56): the universal model is never told that the input is already clean, so it always
re-encodes it. These failures motivate the identity branch of Tasks~2 and~3.

%=====================================================================
\FloatBarrier
\section{Task 2: Classification and Hard-Routed Specialists}
\subsection{Corruption classifier}
The classifier is a VGG-style CNN (stages of two convolution--BatchNorm--ReLU layers and max pooling). The
first stage works at full resolution so that single-pixel salt-and-pepper evidence survives, and the pooled
feature concatenates global average pooling (global blur statistics) with global max pooling (sparse,
extreme evidence such as black rectangles). It is trained with cross-entropy and AdamW~\cite{loshchilov2019adamw}
on balanced batches; model selection uses $1-$macro-F1 on the validation manifest.
\artifacttable{optuna_task2_cls.tex}
The Optuna study (25 trials of 10 epochs; 5 completed, 20 pruned) found a small model: three stages with
16/32/64 channels, dropout 0.18, weight decay $1.6\cdot10^{-3}$, learning rate $3.8\cdot10^{-4}$ and batch 32
(72.8\,k parameters). Dropout, weight decay and the learning rate share the importance (0.40/0.30/0.30),
while the channel plan does not matter, so the task is easy enough for a small network. Most trials were
pruned because the median pruner stops trials that are not better than earlier ones in the same epoch,
which is common when all reasonable configurations reach $>$99\,\% accuracy.
\artifacttable{task2_classifier_per_class.tex}
\begin{figure}[t]\centering\artifactfig{task2_confusion_matrix.png}{confusion matrix}
\caption{Normalised four-class confusion matrix on the test manifest.}\label{fig:cm}\end{figure}
On the 36,690 test records the classifier reaches an accuracy of 99.42\,\%, with macro precision 0.987,
recall 0.994 and F1 0.990. Salt-and-pepper is recognised perfectly. The remaining errors are all
confusions with \emph{clean} (Fig.~\ref{fig:cm}): 0.57\,\% of clean images are predicted as blur and 0.27\,\%
as occlusion, while 0.86\,\% of blurred and 0.80\,\% of occluded images are predicted as clean. All of them
occur at the lowest severity: the accuracy is 97.4\,\% for mild blur ($k=3$, $\sigma=0.7$) and 97.8\,\% for
10\,\% occlusion, and $\geq$99.8\,\% everywhere else. Mildly blurred images of already soft photographs
and small occluders on dark backgrounds are genuinely ambiguous.

\subsection{Specialist autoencoders}
\artifacttable{optuna_task2_spec.tex}
The three specialists share one architecture found by a shared Optuna search, as permitted by the
assignment: every trial trains all three specialists (each only on its own corruption) and minimises their
mean validation $J$. Of 20 trials (11 completed, 9 pruned), the best reduced the mean $J$ from 0.259
(baseline) to 0.174. The selected configuration has a $16\times16\times16$ bottleneck (12$\times$
compression), $C=32$, $\alpha=0.52$ and learning rate $2.0\cdot10^{-3}$ (466\,k parameters per expert); the
learning rate (importance 0.62) and $\alpha$ (0.28) mattered most. The three specialists were then trained
independently for 60 epochs, reaching validation SSIM 0.824 (salt), 0.817 (blur) and 0.715 (occlusion).

\subsection{Hard routing: oracle vs.\ predicted}
\artifacttable{task2_oracle_vs_predicted.tex}
\begin{figure}[t]\centering\artifactfig{task2_oracle_vs_predicted.png}{oracle vs predicted}
\caption{Test SSIM of oracle and predicted routing per condition and severity.}\label{fig:t2routing}\end{figure}
In oracle mode the manifest label selects the expert; in predicted mode the classifier does. A clean
route uses the identity bypass, so the output equals the input exactly (its PSNR is infinite and is capped at
100\,dB in the tables; we therefore compare systems by SSIM). The two modes are almost identical
(Fig.~\ref{fig:t2routing}): the mean SSIM is 0.8018 with oracle and 0.8023 with predicted routing, against
0.672 for the inputs. Predicted routing is even marginally better because some of its ``errors'' are
beneficial: when a mildly blurred image is routed to the identity branch, the output keeps all of its
detail, whereas the blur expert, limited by its bottleneck, lowers the SSIM of such near-clean inputs
(oracle 0.825 vs.\ input 0.944). The comparison of oracle and predicted routing therefore shows that the
classifier is not the bottleneck of the system; the restoration quality of the specialists is.

\subsection{Failure analysis}
\begin{figure}[t]\centering\artifactfig{task2_routing_failures.png}{routing failures}
\caption{Restoration failures caused by classifier errors: clean target, input, oracle route and
predicted route.}\label{fig:t2fail}\end{figure}
Only 214 of 36,690 test records (0.58\,\%) are misrouted. We count a misrouting as a classifier-caused failure
when the predicted route raises $J$ by more than 0.01 compared with the oracle route; this happens 34 times
(Fig.~\ref{fig:t2fail}). Thirty-one of them are clean images sent to the blur (21) or occlusion (10) expert,
which re-encode the image and destroy fine detail; the other three are mildly occluded images sent to the
identity branch, which leaves the black rectangle unrepaired. Conversely, misrouting mild blur or mild
occlusion to the identity branch \emph{lowered} $J$ on average (by 0.091 and 0.068), for the same reason as
above. The most damaging error is therefore a clean image wrongly classified as corrupted.

%=====================================================================
\FloatBarrier
\section{Task 3: Soft Mixture-of-Experts}
\subsection{Architecture and initialisation}
\begin{figure}[t]\centering\artifactfig{arch_moe.pdf}{MoE diagram}
\caption{Soft MoE: gate initialised from the Task 2 classifier, experts from the Task 2 specialists.}\label{fig:moe}\end{figure}
The soft MoE (Fig.~\ref{fig:moe}) computes $w=\mathrm{softmax}(G(\tilde x)/\tau)$ and
$\hat x=w_0\tilde x+w_1A_{salt}(\tilde x)+w_2A_{blur}(\tilde x)+w_3A_{occ}(\tilde x)$. Nothing is randomly
initialised: the gate $G$ is the trained Task~2 classifier (its class order equals the branch order) and the
experts are the three trained specialists (1.47\,M parameters in total). Evaluated without any further
training, this combination already reaches validation SSIM 0.842.

\subsection{Training strategy and joint objective}
$\mathcal{L}_{MoE}=\lambda_1\mathcal{L}_{1}+\lambda_s(1-\mathrm{SSIM})+\lambda_c\mathcal{L}_{CE}+\lambda_b\mathcal{L}_{balance}$,
with the cross-entropy applied to the untempered gate logits, so that $\tau$ only controls how sharply the
gate's belief is turned into weights. Training starts with a warm-up of three epochs in which the experts
are frozen (including their BatchNorm statistics) and only the gate is trained with learning rate $10^{-4}$;
then all parameters are fine-tuned for 15 epochs with the smaller learning rate selected by Optuna. Expert
BatchNorm statistics remain frozen in the joint stage, because every expert now receives all four
conditions, and updating its statistics with inputs it never specialised on would shift its normalisation.

\textbf{Balance regulariser.} The assignment's loss $\mathcal{L}_{balance}=\sum_k(\bar w_k-1/4)^2$ is
computed over balanced batches. For a dense softmax the mean weights sum to one, so their mean is $1/4$ and
Shazeer's importance loss~\cite{shazeer2017moe} $\mathrm{CV}(\bar w)^2=4\sum_k(\bar w_k-1/4)^2$ is the same
loss up to a factor of four (verified by a unit test). As a research-supported alternative we implemented a
marginal-entropy loss $\log4-H(\bar w)$, which is also minimised by uniform usage but whose gradient grows as
a branch dies out. The loss type is itself an Optuna parameter, and the final model was retrained with the
other loss type as an ablation.
\artifacttable{optuna_task3.tex}
\begin{figure}[t]\centering\artifactfig{task3_training_curves.png}{training curves}
\caption{Task 3 training: three warm-up epochs (frozen experts) followed by joint fine-tuning.}\label{fig:t3curves}\end{figure}
The study (20 trials; 15 completed, 5 pruned, either as poor or because of routing collapse) improved the
validation $J$ from 0.093 for the assignment's starting values ($\lambda_1=0.8$, $\lambda_s=0.2$,
$\lambda_c=0.1$, $\lambda_b=0.01$) to 0.082. The selected configuration uses $\tau=2.08$, an equal weighting of
L1 and SSIM ($\lambda_1=0.50$), a small classification weight $\lambda_c=0.010$, a very small balance weight
$\lambda_b=1.9\cdot10^{-4}$ and a joint learning rate of $1.35\cdot10^{-5}$. The reconstruction weighting
(importance 0.52) and the joint learning rate (0.24) mattered most; the balance type had no measurable
influence. In the final run (Fig.~\ref{fig:t3curves}) the warm-up alone improved the validation $J$ from
0.096 to 0.088, and joint fine-tuning reduced it further to 0.080 (SSIM 0.869). Over the same time the
fraction of inputs whose largest weight matches the true label fell from 0.994 to 0.898: the gate learned
to deviate from pure classification where hedging gives a better reconstruction. The balance ablation with
the entropy loss produced the same test SSIM (0.838), which is consistent with the tiny selected $\lambda_b$:
on balanced batches with a well-initialised gate, routing collapse did not occur, and the balance term was
not needed.

\subsection{Routing analysis}
\artifacttable{task3_routing_weights.tex}
\begin{figure}[t]\centering\artifactfig{task3_routing_heatmap.png}{routing heatmap}
\caption{Average routing weight per true condition and severity (test).}\label{fig:heat}\end{figure}
\begin{figure}[t]\centering\artifactfig{task3_weight_distribution.png}{weight distribution}\\
\artifactfig{task3_expert_utilisation.png}{utilisation}
\caption{Routing-weight distributions per true condition and expert utilisation.}\label{fig:util}\end{figure}
\begin{figure}[t]\centering\artifactfig{task3_dominant_examples.png}{dominant}
\caption{Examples in which one branch dominates (largest weight $>0.9$).}\label{fig:dom}\end{figure}
\begin{figure}[t]\centering\artifactfig{task3_mixed_examples.png}{mixed}
\caption{Examples with weights distributed over several branches (largest weight $<0.6$).}\label{fig:mix}\end{figure}
The routing (Table~\ref{tab:task3_routing}, Fig.~\ref{fig:heat}) is sharp for clean (identity 0.92) and
salt-and-pepper inputs (salt expert 0.90--0.99). For blur and occlusion the gate mixes the identity branch
with the matching expert, and the mix depends on severity in a sensible way: the identity weight falls from
0.56 to 0.30 for blur and from 0.53 to 0.26 for occlusion as severity increases, while the expert weight rises
from 0.38 to 0.64 (blur) and from 0.47 to 0.74 (occlusion). For mild corruption the input itself is a good
estimate, so keeping part of it preserves detail that the bottleneck expert would lose. 34.6\,\% of all
test inputs are routed with a largest weight above 0.9 (Fig.~\ref{fig:dom}) and 21.2\,\% with a largest
weight below 0.6 (Fig.~\ref{fig:mix}); the mean routing entropy is 0.54 nats (maximum 1.39). No expert is
inactive: the mean weights are 0.34 (identity), 0.29 (salt), 0.17 (blur) and 0.20 (occlusion), and every
branch is the dominant branch for 21--30\,\% of the inputs (Fig.~\ref{fig:util}). The only weight above 0.3
on an unrelated class is the identity branch on blur and occlusion inputs (0.39 on average), which, as
explained above, is a deliberate hedge rather than a collapse. The salt and blur experts receive at most
0.05 weight on unrelated classes.

\subsection{Results and comparison}
\artifacttable{task3_by_level.tex}
\begin{figure}[t]\centering\artifactfig{task3_system_comparison.png}{system comparison}
\caption{Test SSIM of Task 1, Task 2 (oracle/predicted) and Task 3 per condition and severity.}\label{fig:comp}\end{figure}
The soft MoE is the best restoration system on the test set (overall SSIM 0.838, Fig.~\ref{fig:comp}). It
almost preserves clean images (SSIM 0.998, PSNR 49.5\,dB; not exactly 1, because the gate never gives the
identity exactly all of the weight), and it is the best system at every corruption level. The largest gains
over hard routing are at low severity, where hedging with the identity pays off (mild blur: 0.913 vs.\
0.829; 10\,\% occlusion: 0.876 vs.\ 0.777); at high severity the systems converge (blur 0.761 vs.\ 0.746,
occlusion 0.684 vs.\ 0.656), because then the experts' restoration ability limits all systems.

\subsection{Failure analysis}
The worst Task~3 outputs are the same hard images as in Task~1 (e.g.\ the cat on a black background with
35\,\% occlusion, SSIM 0.723$\to$0.282): the occlusion expert, like the universal model, fills large holes
with a smooth average colour. The four largest degradations are all medium or high occlusions of cats with
dark or highly textured fur, where the inpainted region has the wrong colour. As an additional analysis,
not part of the official protocol, we tested images with \emph{two} corruptions, the situation that
motivates soft routing. For blur+salt-and-pepper at medium severity the SSIM is 0.241 for the input, 0.737
with hard routing and 0.761 with the soft MoE; for occlusion+salt-and-pepper it is 0.252/0.602/0.609 and
for blur+occlusion 0.588/0.668/0.681. The soft MoE is better in every combination, but the gains are small:
because training contained only single corruptions, no expert has learned to remove two corruptions at once.

%=====================================================================
\FloatBarrier
\section{Task 4: Style-Conditioned Face-to-Sketch cGAN}
\subsection{Dataset}
See Section~\ref{sec:fs2k_ref}: the 899/159 training/validation pairs are used for training and model
selection; the official test set (1,046 pairs) only for the final evaluation.

\subsection{cGAN architecture and style conditioning}
\begin{figure}[t]\centering\artifactfig{arch_cgan.pdf}{cGAN diagram}
\caption{U-Net generator and PatchGAN discriminator, each with its own learned style embedding.}\label{fig:cgan}\end{figure}
The generator (Fig.~\ref{fig:cgan}) is a pix2pix U-Net adapted to $128\times128$ inputs: seven stride-2
downsampling stages (to $1\times1$) and a mirrored decoder with skip connections, InstanceNorm~\cite{ulyanov2016instancenorm}
and dropout in the first three decoder blocks (42.0\,M parameters). The discriminator is a $70\times70$
PatchGAN that classifies $14\times14$ overlapping patches of a (photo, sketch) pair (2.8\,M parameters). Each
network has its own learned embedding $e_s\in\mathbb{R}^{16}$ for the three FS2K styles. We compared three ways
of conditioning: spatial replication and concatenation~\cite{choi2018stargan}, conditional normalisation /
FiLM~\cite{devries2017cbn,perez2018film}, and a projection discriminator~\cite{miyato2018projection}. We chose
concatenation because it conditions both networks with very little code and exports directly to ONNX. To
prevent the style signal from fading in deep layers, the generator receives the embedding at the input
(tiled to $128\times128$) and again at the $1\times1$ bottleneck; the discriminator receives it tiled
together with the photo and the sketch. A unit test confirms that changing the style changes the outputs of
both networks.

\subsection{Losses}
$\mathcal{L}_G=\mathrm{BCE}(D(x,G(x,s),s),1)+\lambda_{L1}\,\lVert y-G(x,s)\rVert_1$ and
$\mathcal{L}_D=\tfrac12[\mathrm{BCE}(D(x,y,s),1)+\mathrm{BCE}(D(x,G(x,s),s),0)]$, with binary cross-entropy on
logits, Adam ($\beta_1=0.5$)~\cite{kingma2015adam} and a linear learning-rate decay over the second half of
training, as in pix2pix~\cite{isola2017pix2pix}.

\subsection{Training}
\artifacttable{optuna_task4.tex}
\begin{figure}[t]\centering\artifactfig{task4_gan_losses.png}{GAN losses}
\caption{Discriminator real/fake loss, generator adversarial loss and L1 loss during the final training.}\label{fig:ganloss}\end{figure}
\begin{figure*}[t]\centering\artifactfig[0.9\textwidth]{task4_progression.png}{progression}
\caption{Generator outputs on the same fixed validation photos across the 200 training epochs.}\label{fig:prog}\end{figure*}
Optuna trials used 25-epoch schedules. 22 trials were started: 18 completed, 2 were pruned and 2 were
interrupted when the training process was stopped externally and never finished; they are excluded. The
pix2pix defaults (trial~0: learning rates $2\cdot10^{-4}$, $\lambda_{L1}=100$, dropout 0.5) reached $J=0.327$,
the best trial 0.313. The selected configuration uses $\lambda_{L1}=148.7$, dropout 0.12, a generator learning
rate of $4.5\cdot10^{-4}$, a six times smaller discriminator learning rate of $7.4\cdot10^{-5}$, batch 8, 64 base
channels and a 16-dimensional style embedding. $\lambda_{L1}$ dominated the importance (0.67). Because the
selection objective is pixel-based, it favours a large $\lambda_{L1}$ and therefore conservative, slightly
smoother sketches; this is a known limitation of selecting GANs with paired pixel metrics. The selected
configuration was retrained for the full 200 epochs, and the checkpoint with the best validation $J$
(epoch 113) is deployed.

The losses were logged separately (Fig.~\ref{fig:ganloss}). Training was stable: the discriminator's mean
probability stayed around 0.72 for real and 0.28 for generated pairs from epoch 25 to 200, so neither network
overpowered the other and no mode collapse occurred. The generator's training L1 kept decreasing
(0.284$\to$0.171) while the validation L1 stayed at 0.102--0.103 after epoch 25, so later epochs mainly
improved texture on the training images. Fig.~\ref{fig:prog} shows how the outputs on the same validation
photos develop from blurry outlines into stroke-based sketches.

\subsection{Results}
\artifacttable{task4_test_by_style.tex}
\begin{figure}[t]\centering\artifactfig{task4_style_swap.png}{style swap}
\caption{The same photograph under the three style conditions (test set).}\label{fig:swap}\end{figure}
The assignment-required reconstruction measure is the paired L1 error between the generated and the true
sketch; SSIM and PSNR are reported as additional, research-based measures of structural and pixel
fidelity. We did not use FID: with about 1,000 test images it is strongly biased~\cite{heusel2017fid}, and its
Inception features are not designed for grayscale line drawings. On the official test set the generator
reaches L1 0.111, SSIM 0.462 and PSNR 15.05\,dB (validation: L1 0.102, SSIM 0.482). Style~2 is clearly the
hardest (SSIM 0.372 vs.\ 0.507 for Style~1 and 0.606 for Style~3): its sketches are drawn with dark, dense
strokes whose exact placement varies between drawings, which pixel-wise metrics penalise heavily. Because
the test set contains 619 Style~1, 381 Style~2 and only 46 Style~3 images, the overall mean is dominated by
Styles~1 and~2, and the Style~3 result rests on few images. The style embedding demonstrably controls the
output: generating the same test photo under two different styles changes the sketch by 0.116 (Style 1 vs.\
2), 0.115 (1 vs.\ 3) and 0.097 (2 vs.\ 3) mean absolute intensity, and Fig.~\ref{fig:swap} shows the
visible change from light outlines (Style~1) to dark shaded strokes (Style~2).

\subsection{Failure analysis}
\begin{figure}[t]\centering\artifactfig{task4_style2.png}{style 2 examples}
\caption{Style 2 test examples: best, median and worst by SSIM.}\label{fig:style2}\end{figure}
The worst test results are all Style~2 sketches of \emph{photo1} faces (SSIM 0.13--0.18, Fig.~\ref{fig:style2}).
In these cases the generator places dark hair and shading strokes where the artist drew few lines, or the
other way round; the sketch is plausible but differs from the particular drawing. Because style is confounded
with the photo source in FS2K, the embedding has partly learned the drawing habits of each source, which is
visible when a photo is rendered in a style that never occurred with its source. Finally, at $128\times128$
pixels thin strokes are only one or two pixels wide, so small misalignments already produce large L1 errors.

%=====================================================================
\FloatBarrier
\section{Application Architecture}
\begin{figure}[t]\centering\artifactfig{arch_application.pdf}{application diagram}
\caption{Deployment: React/Tailwind frontend (nginx) and FastAPI + ONNX Runtime backend in Docker Compose.}\label{fig:apparch}\end{figure}
All four systems are served by one web application (Fig.~\ref{fig:apparch}). The React/Tailwind frontend
is built and served by nginx, which forwards \texttt{/api} requests to the backend. The FastAPI backend
validates uploads by content (format, size limit, image dimensions), converts them to RGB and resizes them to
$128\times128$, optionally applies a corruption at runtime with exactly the same code that generated the
training data, runs the ONNX models and returns the images, probabilities, routing weights and timings as
JSON. It provides health-check, universal-restoration, hard-routing, soft-mixture and face-to-sketch
endpoints. The backend contains no PyTorch: it only needs ONNX Runtime, NumPy and Pillow, which a test verifies.
The four workspaces show the input, the restored or generated image, the corruption settings and the
inference time; the Hard-Routed workspace adds the four class probabilities and the selected expert, the
Soft MoE workspace the four routing weights with a stacked bar, and the Face-to-Sketch workspace supports
upload, webcam capture, style selection and download (Figs.~\ref{fig:apphard}--\ref{fig:app}).
\begin{figure}[t]\centering\artifactfig{app_universal.png}{app universal}
\caption{Universal Restoration workspace with the deployed model.}\label{fig:appuni}\end{figure}
\begin{figure}[t]\centering\artifactfig{app_hard.png}{app hard}
\caption{Hard-Routed Restoration workspace: class probabilities and selected expert.}\label{fig:apphard}\end{figure}

\subsection{Interface design in Google Stitch}
The interface was first designed in Google Stitch (Figs.~\ref{fig:stitch1}--\ref{fig:stitch2}): a fixed
sidebar listing the four workspaces and the system status page, a header bar, white rounded cards on a
light canvas, an indigo accent colour, monospace labels for technical values, and a two-column workspace
layout (input and corruption controls on the left; an ``inspection stage'' with tagged image panels,
analysis and telemetry cards on the right). These design tokens and components were transferred to the
React/Tailwind implementation (\texttt{frontend/src/index.css}, \texttt{components/ui.jsx},
\texttt{workspaces/*.jsx}); Fig.~\ref{fig:app} shows the implemented application.
\emph{Note:} the Stitch mock-ups contain placeholder content generated by Stitch (e.g.\ face photographs
in Task~1, hardware names, model names and metric values); these are not results of this project. Only
the visual design was adopted; every value shown by the implemented application comes from the backend.
The main deviations from the mock-ups follow from this: images are shown at the models' $128\times128$
resolution, controls that Stitch invented without a model behind them (e.g.\ a stroke-density slider or SVG
export) were omitted, and the style cards show the three FS2K style categories.
\begin{figure*}[t]\centering
\stitchfig[0.24\textwidth]{stitch_universal.png}\hfill
\stitchfig[0.24\textwidth]{stitch_hard_routing.png}\hfill
\stitchfig[0.24\textwidth]{stitch_soft_moe.png}\hfill
\stitchfig[0.24\textwidth]{stitch_face_to_sketch.png}
\caption{Original Google Stitch designs of the four workspaces: Universal Restoration, Hard-Routed
Restoration, Soft Mixture-of-Experts Restoration and Face-to-Sketch Generator (placeholder content
generated by Stitch).}\label{fig:stitch1}\end{figure*}
\begin{figure}[t]\centering\stitchfig{stitch_system_status.png}
\caption{Google Stitch design of the System status page.}\label{fig:stitch2}\end{figure}
\begin{figure}[t]\centering\artifactfig{app_soft_moe.png}{application screenshot}
\caption{Implemented application (Soft MoE workspace) following the Stitch design.}\label{fig:app}\end{figure}

\FloatBarrier
\section{ONNX Deployment}
\begin{table}[t]\centering\caption{PyTorch vs.\ ONNX Runtime on real validation inputs plus random inputs
(tolerance: \texttt{np.allclose}, atol $10^{-4}$, rtol $10^{-3}$).}\label{tab:onnx}\footnotesize
\setlength{\tabcolsep}{3pt}
\resizebox{\columnwidth}{!}{%
\begin{tabular}{llrrc}\toprule
Model & Output & Max abs.\ diff. & Mean abs.\ diff. & Pass\\\midrule
Task 1 universal DAE & restored & $2.8\cdot10^{-6}$ & $9.7\cdot10^{-8}$ & \checkmark\\
Task 2 classifier & logits & $3.8\cdot10^{-6}$ & $7.6\cdot10^{-7}$ & \checkmark\\
Task 2 salt expert & restored & $2.9\cdot10^{-6}$ & $1.1\cdot10^{-7}$ & \checkmark\\
Task 2 blur expert & restored & $3.5\cdot10^{-6}$ & $1.3\cdot10^{-7}$ & \checkmark\\
Task 2 occlusion expert & restored & $7.6\cdot10^{-6}$ & $1.9\cdot10^{-7}$ & \checkmark\\
Task 3 soft MoE & restored & $3.2\cdot10^{-6}$ & $1.0\cdot10^{-7}$ & \checkmark\\
Task 3 soft MoE & weights & $2.1\cdot10^{-7}$ & $3.4\cdot10^{-8}$ & \checkmark\\
Task 4 generator & sketch & $8.1\cdot10^{-6}$ & $3.1\cdot10^{-7}$ & \checkmark\\\bottomrule
\end{tabular}}\end{table}
All seven deployed models were exported with the \texttt{torch.export}-based ONNX exporter (opset~18) with
a dynamic batch dimension. Each export was validated by running the same inputs (real validation images and
random tensors, at a batch size different from the export batch) through PyTorch and ONNX Runtime and
comparing shapes and values; every output passed, with a maximum absolute difference of $8.1\cdot10^{-6}$,
far below one 8-bit grey level ($3.9\cdot10^{-3}$) (Table~\ref{tab:onnx}). The complete soft MoE (gate,
softmax with the learned $\tau$, three experts and the weighted sum) is a single static graph with two outputs,
so it exports without any architectural change. Hard routing contains data-dependent control flow, so it is
deployed as four graphs, with the argmax and the identity bypass in the backend; this also means that no
expert is executed for inputs predicted as clean. The discriminator is not exported.

\section{Docker and Reproducibility}
The frontend and backend run in two containers started with one command, \texttt{docker compose up --build};
the application is then available at \texttt{http://localhost:8080}. The ONNX models are not stored in the
Git repository but attached to a GitHub release and mounted read-only into the backend container; the backend
has a health check, and the frontend container starts only when the backend is healthy. Ports, the model path
and the upload limit are configurable through environment variables. Reproducibility of the experiments rests
on fixed seeds, the committed split and corruption manifests, the committed Optuna studies, pinned dependency
versions and one driver script (\texttt{scripts/run\_pipeline.py}) that runs every study, training,
evaluation, export and figure step; a SMOKE mode runs the entire pipeline on a tiny subset in a few minutes,
and 67 automated tests cover the corruptions, manifests, models, losses, ONNX consistency and API.

\section{Limitations}
All models work at $128\times128$ pixels, which limits detail in both restoration and sketches. The genuine
bottleneck costs fidelity on clean and mildly corrupted images (Task~1 clean SSIM 0.831); the limited-skip
ablation shows how much is lost. Training contained only one corruption per image, so no system removes
combined corruptions well. Our selection objectives are pixel- and structure-based, which favours smooth
outputs, particularly for the GAN. FS2K style is confounded with the photo source and the Style~3 test set
is small, so the per-style results must be interpreted with care. Bit-exact GPU reproducibility is not
guaranteed by PyTorch, so repeated runs can differ slightly. Finally, the Stitch mock-ups contain
placeholder content and served only as a visual design.

\section{Conclusion}
For blind multi-corruption restoration with a genuine bottleneck, routing to specialists is clearly better
than a single universal model, and soft routing is better than hard routing: the test SSIM increases from
0.786 (universal) to 0.802 (hard) and 0.838 (soft). The decisive advantage of both routed systems is the
identity branch, which keeps clean inputs intact. The hard router's classifier is nearly perfect (99.42\,\%),
so its errors hardly affect the overall quality; the soft MoE's advantage comes from hedging between the
identity and an expert for mild corruptions, a behaviour that emerged from joint training although the gate
was initialised as a classifier. The balance loss turned out to be unnecessary on balanced batches with a
well-initialised gate. For face-to-sketch synthesis, a single pix2pix model with a learned style embedding in
both networks produces three distinguishable styles, with Style~2 the hardest to reproduce. Higher
resolution, training with combined corruptions and perceptual selection criteria are the most promising
next steps.

\appendix[AI Use]
AI assistants were used during this project, and all their output was checked as described here.
\textbf{Claude Code (Anthropic)} was used to extract the requirement checklist from the assignment PDF, to
plan the repository, to implement and debug the data pipeline, models, training, evaluation, ONNX export,
FastAPI backend and React frontend, to write automated tests and documentation, and to find and verify
literature (every reference was checked against its DOI, arXiv or proceedings page). \textbf{Google Stitch}
generated the interface mock-ups. Verification: 67 unit and API tests (e.g.\ corruption ranges, manifest
determinism, SSIM against scikit-image, paired-augmentation alignment, MoE weights summing to one, style
influence on both networks); a SMOKE run of the complete pipeline before the final experiments; inspection of
training curves, logs and generated figures; numerical PyTorch--ONNX comparison of every deployed model; and
manual testing of the application in the browser and in Docker. Problems found and corrected this way included
an incorrect floating-point comparison in a test, a pandas incompatibility in the Task~3 evaluation, an
insufficient Task~1 search space (replaced by a bottleneck-shape search after the first trial) and an
unreliable dependency resolution in the Docker build.

\bibliographystyle{IEEEtran}
\bibliography{references}

\end{document}
